\section{Index of Notation}\label{sec:Notation}

To assist the reader, we summarize the principal symbols, spaces, and functionals used throughout the proof.

\begin{remark}[Notation Conventions]
\textbf{Jang metric:} We use $\bg$ (equivalently written as $\bar{g}$ or $\overline{g}$) for the Jang metric. The macro $\backslash$\texttt{bg} expands to $\overline{g}$.

\textbf{Weight parameters:} On cylindrical ends, $\beta$ denotes the exponential weight in Lockhart--McOwen spaces. On asymptotically flat ends, $\delta$ denotes the polynomial weight. The specific choice $\beta \in (-1, 0)$ for marginally stable MOTS avoids indicial roots.

\textbf{Eigenvalue indexing:} We use \textbf{1-indexing} for eigenvalues of the stability operator $L_\Sigma$. Thus $\lambda_1$ denotes the \textbf{principal (smallest)} eigenvalue, and a stable MOTS satisfies $\lambda_1 \ge 0$. A marginally stable MOTS has $\lambda_1 = 0$, while a strictly stable MOTS has $\lambda_1 > 0$.
\end{remark}

\begin{table}[ht]
\centering
\caption{Metrics, manifolds, and domains.}
\begin{tabular}{l l l}
\hline
\textbf{Symbol} & \textbf{Description} & \textbf{Regularity} \\
\hline
$(M, g, k)$ & Initial data set & Smooth ($C^\infty$) \\
$(\bM, \bg)$ & Jang manifold (graph of $f$) & Lipschitz ($C^{0,1}$) at $\Sigma$ \\
$(\tM, \tg)$ & Conformal deformation ($\tg = \phi^4 \bg$) & $C^0$ at tips $p_k$, Lipschitz at $\Sigma$ \\
$(\tM, \geps)$ & Smoothed manifold (Miao--Piubello) & Smooth ($C^\infty$) \\
$\Sigma$ & Outermost MOTS (horizon) & Smooth embedded surface \\
$\{p_k\}$ & Compactified Jang bubbles & Conical singularities \\
$\mathcal{E}_{cyl}$ & Cylindrical end over $\Sigma$ & $[0,\infty) \times \Sigma$ \\
$\mathcal{E}_{AF}$ & Asymptotically flat end & $M \setminus K$ for compact $K$ \\
\hline
\end{tabular}
\end{table}

\begin{table}[ht]
\centering
\caption{Function spaces and operators.}
\begin{tabular}{l l}
\hline
\textbf{Symbol} & \textbf{Description} \\
\hline
$W^{k,p}_{\delta,\beta}(\bM)$ & Weighted Sobolev space (Lockhart--McOwen) \\
$\Hone$ & $H^1_{\mathrm{loc}}$, locally $H^1$ functions \\
$\Wkp$ & $W^{1,p}_{\mathrm{loc}}$, locally $W^{1,p}$ functions \\
$L_\Sigma$ & Stability operator of MOTS: $L_\Sigma = -\Delta_\Sigma + \tfrac{1}{2}R_\Sigma - |A|^2 - \Ric(\nu,\nu)$ \\
$\Delta_p$ & $p$-Laplacian: $\Delta_p u = \Div(|\nabla u|^{p-2} \nabla u)$ \\
$\mathcal{L}_\phi$ & Lichnerowicz operator: $\mathcal{L}_\phi = -8\Delta + R$ \\
\hline
\end{tabular}
\end{table}

\begin{table}[ht]
\centering
\caption{Key functionals and scalar quantities.}
\begin{tabular}{l l}
\hline
\textbf{Symbol} & \textbf{Description} \\
\hline
$M_{\ADM}(g)$ & ADM mass of metric $g$ \\
$A(\Sigma)$ & Area of surface $\Sigma$ \\
$\mathcal{M}_p(t)$ & AMO monotonicity functional \\
$\mathcal{S}$ & Distributional scalar curvature measure \\
$\Energy_p(u)$ & $p$-energy: $\int |\nabla u|^p \dV$ \\
$\Cap_p(K)$ & $p$-capacity of compact set $K$ \\
$\lambda_1(L_\Sigma)$ & Principal eigenvalue of stability operator \\
\hline
\end{tabular}
\end{table}

\begin{table}[ht]
\centering
\caption{Weight and decay parameters.}
\begin{tabular}{l l l}
\hline
\textbf{Symbol} & \textbf{Range} & \textbf{Role} \\
\hline
$\tau$ & $> 1/2$ & Asymptotic flatness decay rate \\
$\delta$ & $(0, \tau - 1/2)$ & Weight for AF end (order $r^{-\delta}$) \\
$\beta$ & $(-1, 0)$ & Weight for cylindrical ends (order $e^{\beta t}$) \\
$\epsilon$ & $(0, \epsilon_0)$ & Smoothing parameter \\
$\kappa$ & $> 0$ & Surface gravity (blow-up rate: $f \sim \kappa^{-1} \ln s$) \\
\hline
\end{tabular}
\end{table}

\begin{table}[ht]
\centering
\caption{Jump and Interface Quantities.}
\begin{tabular}{l l l}
\hline
\textbf{Symbol} & \textbf{Definition} & \textbf{Context} \\
\hline
$[H]_{\bar{g}}$ & $H_{\text{outer}}^{\bar{g}} - H_{\text{inner}}^{\bar{g}}$ & Mean curvature jump in Jang metric $\bar{g}$ \\
$[H]_{\tilde{g}}$ & Jump in conformal metric $\tilde{g} = \phi^4 \bar{g}$ & Used in Miao smoothing \\
$\tr_\Sigma k$ & Trace of extrinsic curvature & Initial data quantity \\
$\theta^\pm$ & $H_\Sigma \pm \tr_\Sigma k$ & Null expansions \\
\hline
\end{tabular}
\end{table}

\subsection{Detailed Notation Dictionary for Metrics and Jumps}

To avoid ambiguity, we explicitly define the various metrics and jump quantities used in the interface analysis.

\begin{description}
    \item[Initial Data Metric ($g$):] The Riemannian metric on the initial slice $M$.
    \item[Jang Metric ($\bar{g}$):] The metric on the graph $\Sigma \subset M \times \mathbb{R}$, given by $\bar{g} = g + df \otimes df$. It is Lipschitz across the interface $\Sigma$.
    \item[Conformal Jang Metric ($\tilde{g}$):] The metric $\tilde{g} = \phi^4 \bar{g}$ used in the AMO flow. The conformal factor $\phi$ solves the Lichnerowicz equation.
    \item[Smoothed Metric ($g_\epsilon$):] The smooth family of metrics approximating $\tilde{g}$ (or $\bar{g}$) via the Miao--Piubello smoothing.
\end{description}

\textbf{Jump Conventions:}
For a quantity $Q$ discontinuous across a surface $\Sigma$, we define the jump $[Q] = Q^+ - Q^-$, where:
\begin{itemize}
    \item $Q^+$ is the limit from the \textbf{exterior} (asymptotically flat side).
    \item $Q^-$ is the limit from the \textbf{interior} (cylindrical/black hole side).
\end{itemize}
The unit normal $\nu$ points from interior to exterior.

\textbf{Key Identities:}
\begin{itemize}
    \item \textbf{Jang Jump:} $[H]_{\bar{g}} = \tr_\Sigma k$ (under the favorable jump hypothesis).
    \item \textbf{Conformal Jump:} $[H]_{\tilde{g}} = \phi^{-2} [H]_{\bar{g}} + 4\phi^{-3}[\partial_\nu \phi]$.
    \item \textbf{Transmission Condition:} If $[\partial_\nu \phi] = 0$, then $[H]_{\tilde{g}} = \phi^{-2} [H]_{\bar{g}}$.
\end{itemize}

\begin{remark}[Jump Relations]
The mean curvature jumps are related by the conformal factor $\phi$. If the transmission condition $[\partial_\nu \phi] = 0$ holds, then:
\[
[H]_{\tilde{g}} = \phi^{-2} [H]_{\bar{g}}.
\]
The fundamental identity relating the Jang metric jump to initial data is $[H]_{\bar{g}} = \tr_\Sigma k$ (Lemma~\ref{lem:TrappedMeanCurvatureJump}).
\end{remark}

\begin{table}[ht]
\centering
\caption{Key theorems and their roles.}
\small
\begin{tabular}{@{}lll@{}}
\hline
\textbf{Theorem} & \textbf{Content} & \textbf{Section} \\
\hline
\ref{thm:penroseinitial} & Primary result (outermost MOTS) & \S\ref{sec:intro} \\
\ref{thm:MainTheorem} & Conditional Penrose inequality & \S\ref{sec:penrose_conjecture} \\
\ref{thm:CompleteProof} & Consolidated proof & \S\ref{sec:Consolidated} \\
\ref{thm:MaxAreaTrapped} & Maximum area trapped surface & \S\ref{sec:Overview} \\
\ref{thm:AMOMonotonicity} & AMO level set monotonicity & \S\ref{sec:AMO} \\
\ref{thm:CompleteMeanCurvatureJump} & Mean curvature jump & \S\ref{sec:Jang} \\
\hline
\end{tabular}
\end{table}

\begin{table}[ht]
\centering
\caption{Sign conventions summary.}
\begin{tabular}{l l}
\hline
\textbf{Quantity} & \textbf{Convention} \\
\hline
Mean curvature $H$ & $H > 0$ for outward-bending surfaces \\
Null expansion $\theta^\pm$ & $\theta^\pm = H \pm \tr_\Sigma k$ \\
Trapped surface & $\theta^+ \le 0$ (outer trapped) \\
MOTS & $\theta^+ = 0$ (marginally outer trapped) \\
Stability operator $L_\Sigma$ & $\lambda_1 \ge 0$ for stable MOTS \\
Mean curvature jump $[H]$ & $[H] = H^+ - H^-$ (exterior minus interior) \\
Scalar curvature $R$ & Round sphere has $R > 0$ \\
\hline
\end{tabular}
\end{table}

\begin{table}[ht]
\centering
\caption{Geometric Quantities and Jumps.}
\begin{tabular}{l l l}
\hline
\textbf{Symbol} & \textbf{Description} & \textbf{Sign Convention} \\
\hline
$\theta^+$ & Outer null expansion & $\le 0$ for trapped \\
$\theta^-$ & Inner null expansion & $< 0$ for trapped \\
$\tr_\Sigma k$ & Trace of extrinsic curvature on $\Sigma$ & Favorable if $\ge 0$ \\
$[H]_{\bar{g}}$ & Mean curvature jump in Jang metric & $[H]_{\bar{g}} = \tr_\Sigma k$ \\
$[H]_{\tilde{g}}$ & Jump in conformal metric & $[H]_{\tilde{g}} = \phi^{-2}[H]_{\bar{g}}$ \\
$\lambda_1$ & Principal eigenvalue of stability op. & Stable if $\ge 0$ \\
\hline
\end{tabular}
\end{table}

